\documentclass[12pt, a4paper]{article}
\usepackage{../../../myconfig} % Gọi file cấu hình từ ngoài gốc vào

% ==========================================
% BẮT ĐẦU NỘI DUNG TÀI LIỆU
% ==========================================
\begin{document}

% Truyền 3 tham số: {Nhãn cam} {Chữ to chính giữa} {Chữ ở góc phải Header}
\titlebox{Chủ đề: Hàm số}{Hàm doanh thu, lợi nhuận, tối ưu chi phí}{Hàm số: Doanh thu lợi nhuận chi phí}

\drawdashlines{20}
\newpage
\drawdashlines{25}


\questionShortAnswer{20}{ví dụ}{Sau khi kết thúc kì thi THPTQG 2027 thì lứa Trâu đậu Đại học Trăm Mâm, nhưng sau khi tính toán thì bọn Trâu đã tính rằng nếu mời x khách thì tổng tiền mừng nhận được sẽ là 
\( F(x) = -0,03x^2 + 800x \)
(nghìn đồng). Trong khi đó chi phí ăn bình quân cho mỗi người là      
\(       G(x) = \dfrac{50000}{x} + 450 \)  
(nghìn đồng). Giả sử mời người nào cũng đi. Hãy tính thử phải mời bao nhiêu người để được lợi nhuận lớn nhất}


% --- CÂU 11 ---
\questionShortAnswer{20}{1}{Một công ty sử dụng xe tải để vận chuyển các sản phẩm của họ. Để tính toán chi phí, quản lý của công ty này đã mô hình hóa lượng xăng tiêu thụ ước tính theo hàm số \( G(x) = \dfrac{1}{100}\left(\dfrac{3000}{x} + x\right) \) (lít/km), với \( x \) (km/h) là vận tốc của xe tải (giả định rằng vận tốc này duy trì không đổi trong suốt quá trình di chuyển). Biết rằng giá xăng là 20 nghìn đồng/lít. Ngoài ra, công ty cần thuê người lái xe tải với mức lương 500 nghìn đồng mỗi giờ để lái xe tải cho công ty này, và quãng đường mà người lái xe di chuyển để vận chuyển các sản phẩm của công ty dài 400km. Hỏi người lái xe tải phải duy trì vận tốc của xe bằng bao nhiêu km/h để tổng chi phí của công ty là thấp nhất (kết quả làm tròn đến hàng phần mười)?}

% --- CÂU 12 ---
\questionShortAnswer{20}{2}{Ở cửa hàng thời trang TRÂU Shop, anh Huy có một sản phẩm thời trang đặc biệt mà giá bán của nó thay đổi theo thời gian. Sau \( t \) tháng kể từ khi ra mắt sản phẩm, giá bán của sản phẩm đó thay đổi theo công thức  
\(p(t) = \dfrac{2t\sqrt{t}}{25} - t + 170 \ (\text{nghìn đồng}).\)
Tuy nhiên, nhu cầu mua sản phẩm này cũng thay đổi theo giá bán: Nếu giá bán của sản phẩm là \( p \) (nghìn đồng), dự tính mỗi tháng anh Huy sẽ bán được  
\(D(p) = \dfrac{40000}{p}\)
sản phẩm. Hỏi anh Huy có thể bán được tối đa bao nhiêu sản phẩm thời trang đặc biệt trên trong một tháng?}


% --- CÂU 14 ---
\questionShortAnswer{20}{3}{Một cửa hàng phân phối gạo với chi phí mua vào là 25 nghìn đồng/kg và bán ra với giá 30 nghìn đồng/kg. Với giá bán này thì số gạo bán được trong một tháng là 10000kg. Để đẩy mạnh hơn nữa doanh số tiêu thụ gạo trong một tháng, cửa hàng dự định giảm giá bán và ước tính rằng nếu giảm 1 nghìn đồng/kg thì số lượng gạo bán ra trong một tháng sẽ tăng thêm 5000kg. Cửa hàng phải định giá bán gạo mới là bao nhiêu nghìn đồng một kilôgam thì lợi nhuận thu được trong tháng cao nhất?}


% --- CÂU 16 ---
\questionShortAnswer{20}{4}{Bác Thúy có một bãi nuôi cá, mỗi ngày thu hoạch được 1 tấn cá. Nếu bác bán cho thương lái với giá 30 nghìn đồng/kg thì hết cá, nếu giá bán cứ tăng thêm 5 nghìn đồng/kg thì số cá thừa sẽ tăng thêm 20kg. Số cá thừa này được bán để làm thức ăn cho động vật với giá 10 nghìn đồng/kg. Hỏi doanh thu lớn nhất từ việc bán cá mỗi ngày của bác Thúy là bao nhiêu triệu đồng?}

% --- CÂU 17 ---
\questionShortAnswer{20}{5}{Một công ty sản xuất dụng cụ thể thao nhận được một đơn đặt hàng sản xuất 10000 quả bóng pickleball trong 1 năm. Công ty này sở hữu một số máy móc, mỗi máy có thể sản xuất 50 quả bóng trong một giờ. Chi phí thiết lập các máy này là 250 nghìn đồng cho mỗi máy. Khi được thiết lập, hoạt động sản xuất sẽ hoàn toàn diễn ra tự động dưới sự giám sát. Số tiền phải trả cho người giám sát là 200 nghìn đồng một giờ. Số máy móc công ty nên sử dụng là bao nhiêu để chi phí hoạt động là thấp nhất?}

% --- CÂU 18 ---
\questionShortAnswer{20}{6}{Để hạn chế vi phạm thời gian làm việc đối với các công nhân, giám đốc công ty quyết định xử lý bằng cách phạt tiền. Nhờ sự giám sát chặt chẽ của các quản đốc, giám đốc công ty biết được rằng trong một tháng, tỉ lệ công nhân vi phạm đúng \( k \) lần (\( k = 1 \) hoặc \( k = 2 \)) là \( t_k = \dfrac{N_k}{N} \) (trong đó \( N_k \) là số công nhân vi phạm đúng \( k \) lần, \( N \) là tổng số công nhân) và mức phạt cho mỗi lần vi phạm có mối liên hệ như sau: Nếu mỗi công nhân nộp phạt \( x \) nghìn đồng (\( 50 \leq x \leq 500 \)) khi vi phạm lần thứ nhất và nộp phạt \( x + 50 \) nghìn đồng khi vi phạm lần thứ hai thì tỉ lệ công nhân vi phạm lần lượt là \( t_1 = \dfrac{64}{x + 15} \) và \( t_2 = \dfrac{8}{x - 25} \) (giả sử rằng không có công nhân nào vi phạm quá 2 lần trong một tháng). Biết rằng tổng số công nhân bằng 2400, khi đó tổng số tiền nộp phạt của các công nhân vi phạm trong một tháng ít nhất là bao nhiêu triệu đồng (kết quả làm tròn đến hàng đơn vị)?}



% --- CÂU 3 ---
\questionTrueFalse{15}{7}{Chị Dung là chủ của thương hiệu Lowlands Coffee. Thương hiệu này tập trung vào xuất khẩu cà phê sang nước ngoài. Người ta ước tính rằng, nếu giá bán của mỗi kilogram cà phê ở nước ngoài là \( x \) (nghìn đồng) thì mỗi tuần khách hàng sẽ mua khoảng \( D(x) = \dfrac{17500}{x^2} \) kilogram cà phê, với \( D(x) \) là hàm số biểu thị nhu cầu mua cà phê của khách hàng. Ngoài ra, chị Dung còn khảo sát thấy rằng giá bán cà phê sau \( t \) tuần tính từ khi bắt đầu tiếp cận thị trường nước ngoài cũng thay đổi theo hàm số \( x(t) = 0,2t^2 - t + 60 \) (đơn vị: nghìn đồng).}{
    \textbf{a)} & Doanh thu mà thương hiệu của chị Dung đạt được sau \( t \) tuần tính từ khi bắt đầu tiếp cận thị trường nước ngoài là \( R(t) = \dfrac{17500}{(0,2t^2 - t + 60)^2} \) (đơn vị: nghìn đồng).
    & \textcolor{amberorange}{$\bigcirc$} & \textcolor{amberorange}{$\bigcirc$} \\ \hline
    
    \textbf{b)} & Giá bán cà phê đạt mức thấp nhất sau 2,5 tuần tính từ khi bắt đầu tiếp cận thị trường nước ngoài.
    & \textcolor{amberorange}{$\bigcirc$} & \textcolor{amberorange}{$\bigcirc$} \\ \hline
    
    \textbf{c)} & Khi giá bán cà phê theo thời gian đạt mức thấp nhất thì thương hiệu của chị Dung cũng đạt mức doanh thu cao nhất.
    & \textcolor{amberorange}{$\bigcirc$} & \textcolor{amberorange}{$\bigcirc$} \\ \hline
    
    \textbf{d)} & Sau thời gian dài (có thể coi như tiến đến vô cùng), nhu cầu mua cà phê cũng như doanh thu của thương hiệu cà phê sẽ tiến về 0.
    & \textcolor{amberorange}{$\bigcirc$} & \textcolor{amberorange}{$\bigcirc$} \\
}

% --- CÂU 4 ---
\questionTrueFalse{15}{8}{Bác Thương có một nông trại nuôi tôm, mỗi ngày bác thu hoạch được nửa tấn tôm. Nếu bác bán cho thương lái với giá 120 nghìn đồng/kg tôm thì bác bán được hết số tôm. Nếu giá bán cứ tăng thêm 1 nghìn đồng/kg thì bác sẽ bị thừa 5kg tôm. Để tránh lãng phí, bác Thương quyết định sẽ bán số tôm thừa này cho các hộ chăn nuôi làm thức ăn cho động vật với giá là 15 nghìn đồng/kg. Xét \( R(x) \) là hàm số biểu thị doanh thu từ việc bán tôm của bác Thương, với \( x \) là số kilogram tôm mà bác đã bán cho thương lái \( (0 \leq x \leq 500) \).}{
    \textbf{a)} & Nếu bác bán cho thương lái được 400kg tôm thì tổng số tiền bác thu được từ việc bán tôm cho các hộ chăn nuôi là 1,5 triệu đồng.
    & \textcolor{amberorange}{$\bigcirc$} & \textcolor{amberorange}{$\bigcirc$} \\ \hline
    
    \textbf{b)} & Hàm số biểu thị doanh thu từ việc bán tôm cho các hộ chăn nuôi là \( T(a) = 10a \) (nghìn đồng), với \( a \) là số kilogram tôm mà bác đã bán cho các hộ chăn nuôi.
    & \textcolor{amberorange}{$\bigcirc$} & \textcolor{amberorange}{$\bigcirc$} \\ \hline
    
    \textbf{c)} & Biết rằng phương trình \( R(x) = 0 \) có hai nghiệm phân biệt, khi đó tích của hai nghiệm này là một số chính phương.
    & \textcolor{amberorange}{$\bigcirc$} & \textcolor{amberorange}{$\bigcirc$} \\ \hline
    
    \textbf{d)} & Doanh thu tối đa từ việc bán tôm của bác Thương lớn hơn 60 triệu đồng.
    & \textcolor{amberorange}{$\bigcirc$} & \textcolor{amberorange}{$\bigcirc$} \\
}
\end{document}